% Generated draft: requires mathematical and source review.
\documentclass[11pt]{article}
\usepackage[margin=25mm]{geometry}
\usepackage{fontspec}
\setmainfont{Latin Modern Roman}
\usepackage{amsmath,amssymb,mathtools,unicode-math}
\setmathfont{Latin Modern Math}
\usepackage{enumitem,needspace,xurl,hyperref}
\setlength{\parindent}{0pt}
\setlength{\parskip}{5pt}
\setlength{\emergencystretch}{3em}
\newcommand{\N}{\mathbb{N}}
\newcommand{\Z}{\mathbb{Z}}
\newcommand{\Q}{\mathbb{Q}}
\newcommand{\R}{\mathbb{R}}
\newcommand{\C}{\mathbb{C}}
\newcommand{\F}{\mathbb{F}}
\newcommand{\A}{\mathbb{A}}
\newcommand{\PP}{\mathbb{P}}
\newcommand{\E}{\mathbb{E}}
\newcommand{\eps}{\varepsilon}
\newcommand{\dd}{\,\mathrm{d}}
\begin{document}
\section*{UnsolvedMath ID: 701625}
\textbf{Source problem code:} catalogue entry
\subsection{Definitions and mathematical statement}
Let \(F\) be a free group of rank \(r\geq 2\), and let \(n\) be a positive integer. Write \(F^n\) for the set of ordered \(n\)-tuples of elements of \(F\), and \(\operatorname{Aut}(F)\) for its automorphism group. Each \(\phi\in\operatorname{Aut}(F)\) acts diagonally by
\[
\phi^{(n)}(w_1,\ldots,w_n)=(\phi(w_1),\ldots,\phi(w_n)),
\qquad
\phi^{(n)}(S)=\{\phi^{(n)}(\mathbf w):\mathbf w\in S\}.
\]
\textbf{Stable Small Cancellation Conjecture (P. E. Schupp).} For every such \(F\) and every positive integer \(n\), there exists a generic subset \(S\subseteq F^n\) such that, for every \(\phi\in\operatorname{Aut}(F)\), the set \(\phi^{(n)}(S)\) satisfies \(C'(1/6)\). The same subset \(S\) must work simultaneously for all automorphisms.

\textbf{Qualifications needed for a precise formulation.} The supplied record does not define either the meaning of ``generic'' or how \(C'(1/6)\) is applied to a subset of \(F^n\). Thus the preceding statement preserves the source wording but does not specify a complete mathematical predicate.

For context, a usual word-theoretic definition of \(C'(1/6)\) begins with a chosen free basis \(A\) and a symmetrized collection \(\mathcal R\) of nonempty cyclically reduced words over \(A\cup A^{-1}\); symmetrized means closed under inverses and cyclic permutations. A piece is a nonempty word that is an initial segment of two distinct words in \(\mathcal R\). The condition is
\[
\forall R\in\mathcal R\;\forall P\,
\bigl(P\text{ is a piece and an initial segment of }R
\ \Longrightarrow\ |P|_A<\tfrac16|R|_A\bigr).
\]
To use this definition here, one must specify the basis, the treatment of the images of tuple entries as relators, and their reduction and symmetrization. One possible reading requires this condition separately for the relators associated with every tuple in \(\phi^{(n)}(S)\). A collective reading would instead compare relators from different tuples as well. The supplied record does not distinguish these readings, and neither is silently selected here. Likewise, no counting model, length convention, or limiting density defining genericity is supplied.
\subsection{Short English statement}
Schupp's conjecture asks whether, for every free group of rank at least two and every positive integer \(n\), one can choose a generic set of \(n\)-tuples whose images satisfy \(C'(1/6)\) under every automorphism. The same set must work for all automorphisms. The supplied record leaves genericity and the application of small cancellation to a set of tuples unspecified.
\subsection{Sources}
\begin{itemize}
\item Local UnsolvedMath registry, permanent ID \(701625\), exact source problem code \texttt{catalogue entry}, title \emph{Kourovka Notebook Problem 19.77}. Supplied upstream revision: \texttt{d371ee917ac9647fbb06cd6fc529deec6f3dda82}.
\item \emph{The Kourovka Notebook (Unsolved Problems in Group Theory)}, Problem 19.77, attributed in the supplied record to the 19th issue (2018) and the 21st issue edition (2026), arXiv:1401.0300v48, \url{https://arxiv.org/abs/1401.0300v48}. The supplied collection description names E. I. Khukhro and V. D. Mazurov as editors.
\item Proposed by P. E. Schupp; supplied proposal year: 2018. The local record labels the problem \emph{open}. This is supplied status metadata, not an independently verified assessment. The record was updated at \texttt{2026-10-16T14:00:00Z}; that timestamp is not identified as a literature-verification date.
\end{itemize}
\end{document}
